\begin{figure}[p]
\centering
\includegraphics[width=\textwidth]{supplementary/curated/branch_conflict_controls.pdf}
\caption{\textbf{Branch-selective credit limits interference and preserves update direction.} \textbf{A}, Two-branch learning in ten paired seeds. Points show seeds; large symbols and bars are means and 95\% paired-seed bootstrap intervals. Exact path and the equivalent gated-point rule supply branch selection; the context-0 forgetting of the same cohort is drawn at the main-text type scale in main Fig.~2E. \textbf{B--D}, Shared-versus-exact gradient cosine in the separate branch-conflict experiment with two, four or eight branches. Solid curves evaluate common exact-learning states at epochs 0, 50 and 250; dashed curves evaluate the shared learner's own final state. Curves and shading use twenty paired seeds and 95\% bootstrap intervals. An exact quadratic interference calculation is given in the accompanying derivation. Learned-state geometry can depart from the initialization mean-field approximation.}
\label{fig:si_branch_conflict_controls}
\label{fig:supp_prospective_diagnostics}
\label{fig:supp_path_demand}
\end{figure}

\begin{figure}[p]
\centering
\includegraphics[width=\textwidth]{supplementary/curated/ancestry_coefficients.pdf}
\caption{\textbf{Available route span and learned route coefficients are different constraints.} \textbf{A}, Trained field capture by route budget in the eight-context factorial. Ancestry, depth and random partitions have coincident capture at matched bandwidth; derangement removes the target coordinate. Capture therefore does not explain the small ancestry-specific advantage, and capture alone does not predict trained utility: the within-seed Spearman correlation between capture and utility is about 0.82 over all budgets and 0.64 for $K<8$, and is retained numerically in Source Data. \textbf{B}, Held-out accuracy versus calibration set size (16, 64 and 256 examples) at cue-noise standard deviation 0.5 and zero delay; 30 epochs are fixed, so both data and computation increase. \textbf{C}, Cue delay at 256 calibration trials and noise SD 0.5; delays mismatch independent contexts and are measured in trials, not physiological time. In \textbf{B,C} the dashed curves marked with an asterisk use maximum-probability route selection by the same frozen encoder, an exploratory paired sensitivity specified after the soft outcomes with no encoder refit. Curves, diamonds and bars are means and 95\% paired-seed bootstrap intervals over twenty fresh analysis seeds. The cue-noise sensitivity of the same encoder is drawn at the main-text type scale in main Fig.~3F. The calibration cue and activation targets are supplied resources. No serial forward dendritic computation is present.}
\label{fig:si_ancestry_coefficients}
\label{fig:supp_partition_residual}
\label{fig:supp_local_context_encoder}
\end{figure}
